% Equivalente electrico del Taller 6 (analogia esfuerzo-voltaje).
% Se compila aparte:  pdflatex circuito_equivalente.tex
\documentclass[border=4pt]{standalone}
\usepackage[T1]{fontenc}
\usepackage{lmodern}
\usepackage{amsmath}
\usepackage[american, cute inductors]{circuitikz}
\usetikzlibrary{calc}
\ctikzset{bipoles/length=0.85cm, resistors/scale=0.8, capacitors/scale=0.8,
          inductors/scale=0.8, sources/scale=0.8, diodes/scale=0.7,
          batteries/scale=0.8, switches/scale=0.75}
\definecolor{flujo}{RGB}{30,80,220}
\definecolor{esfuerzo}{RGB}{200,20,20}
\definecolor{grupo}{RGB}{120,120,120}
% flujo horizontal (azul) sobre el cable
\newcommand{\fl}[2]{\draw[->, flujo, thick] ($(#1)+(0,0.3)$) -- ++(0.55,0)
                    node[midway, above, font=\small]{#2};}
% flujo vertical hacia abajo (azul)
\newcommand{\fld}[2]{\draw[->, flujo, thick] (#1) -- ++(0,-0.55)
                     node[midway, right, font=\small]{#2};}
% esfuerzo vertical (rojo) entre los cables, etiqueta a la izquierda
\newcommand{\ef}[3]{\draw[->, esfuerzo, thick] (#1) -- ++(0,#2)
                    node[midway, left, font=\small]{#3};}
% esfuerzo horizontal (rojo) hacia el terminal positivo
\newcommand{\efh}[3]{\draw[->, esfuerzo, thick] (#1) -- (#2)
                     node[pos=0.5, below, font=\small]{#3};}
\tikzset{caja/.style={draw, fill=gray!10, align=center},
         nota/.style={font=\small, text=grupo},
         malla/.style={font=\small\bfseries, text=black!75},
         grupo/.style={draw=grupo, dashed, rounded corners=3pt}}
% amortiguador de tope: 1/K y gamma en serie, puenteados por un interruptor
% cerrado que se abre durante el contacto.  #1 = x inicial, #2 = y del cable,
% #3 = K, #4 = gamma, #5 = condicion de contacto, #6 = nombre del tope
\newcommand{\tope}[6]{%
  \draw (#1,#2) to[C, l_=$1/#3$] ({#1+1.3},#2) to[R, l_=$#4$] ({#1+2.6},#2);
  \draw (#1,#2) -- (#1,{#2+1.0}) to[ncs, l=$#5$] ({#1+2.6},{#2+1.0}) -- ({#1+2.6},#2);
  \node[nota] at ({#1+1.3},{#2+1.85}) {#6};}
\begin{document}
\begin{circuitikz}[font=\small, line width=0.6pt]

% =================================================== MALLA 1 y 2: electrica
% fuente monofasica y convertidor
\draw (-0.5,1.0) to[sV, l=$V_s(t)$] (-0.5,2.6);
\draw (-0.5,2.6) -- (0.55,2.6) (-0.5,1.0) -- (0.55,1.0);
\node[caja, minimum width=2.05cm, minimum height=3.7cm] (conv) at (1.575,1.85)
     {\textbf{AC--AC}\\$1\phi\to3\phi$\\[5pt]razón $a$\\[5pt]{\scriptsize $v_\alpha = V_s(t)$}\\{\scriptsize $v_\beta = V_s(t-T/4)$}};
% tres fases en estrella (salida del convertidor) con R_s y diodo de catodo comun
\draw (3.0,0) -- (3.0,3.4);
\foreach \y/\f in {3.4/a, 2.15/b, 0.9/c}{
  \draw (3.0,\y) to[sV, l=$v_\f$] (4.5,\y) to[R, l_=$R_s$] (5.9,\y)
        to[D*, l=$D_\f$] (7.3,\y) -- (7.8,\y);
  \fl{4.95,\y}{$i_{D\f}$}
}
\draw[grupo] (2.8,0.4) rectangle (4.75,4.15);
\node[nota, align=center] at (3.78,4.45) {salida del convertidor};
\draw[->, dashed, grupo] (2.6,2.15) -- (2.8,2.15);
\node at (3.0,0) [below left, font=\small] {$n$};
% bus de CC: catodos unidos y condensador
\draw (7.8,0.9) -- (7.8,3.4) -- (9.0,3.4);
\draw (3.0,0) -- (13.4,0);
\draw (9.0,3.4) to[C, l_=$C$] (9.0,0);
\fl{8.0,3.4}{$i_D$}
\ef{9.55,0.4}{2.6}{}
\node[esfuerzo, right] at (9.55,1.7) {$U$};
\node[nota] at (5.4,-0.45) {rectificador de 3 diodos};
% armadura y girador
\draw (9.0,3.4) -- (10.0,3.4) to[R, l_=$R_a$] (11.4,3.4) to[L, l_=$L_a$] (12.8,3.4) -- (13.4,3.4);
\fl{10.1,3.4}{$i_a$}
\ef{13.15,0.4}{2.6}{$u_e$}
\node[caja, minimum width=1.2cm, minimum height=3.4cm] at (14.0,1.7) {\textbf{GY}\\$K_a$};
\draw (14.6,3.4) -- (16.2,3.4) node[ocirc]{} node[right]{$a$};
\draw (14.6,0) -- (16.2,0) node[ocirc]{} node[right]{$b$};
\fl{14.8,3.4}{$\omega_m$}
\ef{15.8,0.4}{2.6}{$M_m$}
\node[nota] at (14.0,-0.45) {motor};
\node[malla] at (4.6,-1.05) {Malla 1: fuente, convertidor, rectificador y bus de CC};
\node[malla] at (12.6,-1.05) {Malla 2: armadura};

% =================================================== MALLAS 3 y 4: rotacional (Taller 3)
\begin{scope}[yshift=-5.6cm]
\draw (0,3.4) node[ocirc]{} node[left]{$a$} -- (0.3,3.4) to[L, l_=$J_1$] (1.7,3.4)
      to[R, l_=$\beta_1$] (3.1,3.4) -- (4.9,3.4);
\draw (0,0) node[ocirc]{} node[left]{$b$} -- (4.9,0);
\draw (3.8,3.4) to[C, l_=$K_1$] (3.8,0);
\fl{0.35,3.4}{$\omega_m$}
\fld{4.05,2.8}{$\omega_{k1}$}
\node[caja, minimum width=1.2cm, minimum height=3.4cm] at (5.5,1.7) {\textbf{T}\\$N$};
\node[nota] at (5.5,-0.45) {engranajes};
\draw (6.1,3.4) -- (6.4,3.4) to[L, l_=$J_2$] (7.8,3.4) to[R, l_=$\beta_2$] (9.2,3.4) -- (10.9,3.4);
\draw (6.1,0) -- (14.2,0);
\fl{6.1,3.4}{$\omega_2$}
\draw (9.9,3.4) to[C, l_=$K_2$] (9.9,0);
\fld{10.15,2.8}{$\omega_{k2}$}
\draw (10.9,3.4) to[L, l_=$J_p$] (12.3,3.4) -- (14.2,3.4);
\fl{12.6,3.4}{$\omega_p$}
\ef{13.85,0.4}{2.6}{$M_r$}
\node[caja, minimum width=1.2cm, minimum height=3.4cm] at (14.8,1.7) {\textbf{T}\\$r$};
\node[nota] at (14.8,-0.45) {rueda};
\draw (15.4,3.4) -- (16.6,3.4) node[ocirc]{} node[right]{$c$};
\draw (15.4,0) -- (16.6,0) node[ocirc]{} node[right]{$d$};
\fl{15.5,3.4}{$v_A$}
\ef{16.25,0.4}{2.6}{$F_r$}
\node[malla] at (2.4,-1.05) {Malla 3: eje del motor};
\node[malla] at (10.2,-1.05) {Malla 4: eje secundario y rueda};
\end{scope}

% =================================================== MALLA 5a: cabina A
\begin{scope}[yshift=-11.6cm]
\draw (0,2.4) node[ocirc]{} node[left]{$c$} -- (0.4,2.4)
      to[L, l=$m_A$] (1.8,2.4)
      to[battery1, l=$m_A g\,\mathrm{sen}\,\theta$, invert] (3.4,2.4)
      to[R, l=$\gamma$] (4.8,2.4) -- (5.6,2.4);
\node[esfuerzo] at (2.25,2.05) {$+$}; \node[esfuerzo] at (2.95,2.05) {$-$};
\tope{5.6}{2.4}{K_S}{\gamma_S}{\text{abre si } x_A>x_S}{tope superior de A}
\draw (8.2,2.4) -- (8.8,2.4);
\tope{8.8}{2.4}{K_B}{\gamma_B}{\text{abre si } x_A<0}{tope inferior de A}
\draw (11.4,2.4) -- (13.6,2.4) node[ocirc]{} node[right]{$e$};
\draw (0,0) node[ocirc]{} node[left]{$d$} -- (13.6,0) node[ocirc]{} node[right]{$f$};
\fl{11.75,2.4}{$v_A$}
\efh{11.4,1.2}{0.4,1.2}{$F_A$: tensión del cable A}
\ef{13.25,0.3}{1.8}{$-F_B$}
\draw[grupo] (0.25,0.75) rectangle (11.65,4.65);
\node[malla, anchor=west] at (0.4,4.35) {Cabina A (sube)};
\node[malla] at (6.8,-0.5) {Malla 5a: cabina A, recorrida por $v_A$};
\end{scope}

% =================================================== MALLA 5b: cabina B
\begin{scope}[yshift=-17.0cm]
\draw (0,2.4) node[ocirc]{} node[left]{$e$} -- (0.4,2.4)
      to[L, l=$m_B$] (1.8,2.4)
      to[battery1, l=$m_B g\,\mathrm{sen}\,\theta$] (3.4,2.4)
      to[R, l=$\gamma$] (4.8,2.4) -- (5.6,2.4);
\node[esfuerzo] at (2.25,2.05) {$-$}; \node[esfuerzo] at (2.95,2.05) {$+$};
\tope{5.6}{2.4}{K_B}{\gamma_B}{\text{abre si } x_A>x_S}{tope inferior de B}
\draw (8.2,2.4) -- (8.8,2.4);
\tope{8.8}{2.4}{K_S}{\gamma_S}{\text{abre si } x_A<0}{tope superior de B}
\draw (11.4,2.4) -- (12.6,2.4) -- (12.6,0);
\draw (0,0) node[ocirc]{} node[left]{$f$} -- (12.6,0);
\fld{12.85,1.6}{$v_A=-v_B$}
\efh{11.4,1.2}{0.4,1.2}{$-F_B$: la cabina B entrega esfuerzo al lazo}
\draw[grupo] (0.25,0.75) rectangle (11.65,4.65);
\node[malla, anchor=west] at (0.4,4.35) {Cabina B (baja)};
\node[nota, align=left, anchor=west] at (14.6,2.2) {peso de B con\\polaridad invertida:\\$F_r = F_A - F_B$};
\node[malla] at (6.8,-0.5) {Malla 5b: cabina B, recorrida por la misma corriente $v_A$};
\end{scope}
\end{circuitikz}
\end{document}
